\begin{Theorem}
\label{thm:hilbert-basis-equivalent-characterizations}
设 $\{x_i\}_{i\in I}$ 为 Hilbert 空间 $H$ 中的标准正交集. 那么以下条件等价:\begin{enumerate}[label=(\roman*)]
\item $\{x_i\}_{i\in I}$ 是一组基.
\item 对每个 $x,y\in H$, 有
\begin{equation}\label{eq:hilbert-basis-inner-product-expansion}
(x,y)=\sum_{i\in I}(x,x_i)(x_i,y).
\end{equation}

\item 对每个 $x\in X$, 有
\begin{equation}\label{eq:hilbert-basis-parseval-characterization}
\|x\|^2=\sum_{i\in I}|(x_i,x)|^2.
\end{equation}

\item 集合 $\{x_i\}_{i\in I}$ 极大, 即不存在以它为真子集的标准正交集. 换言之, 若 $x$ 与每个 $x_i$ 正交, 则 $x=0$.\end{enumerate}
\end{Theorem}

\begin{Proof}
(i)$\Rightarrow$(ii): 有
\begin{equation}\label{eq:hilbert-basis-double-sum-inner-product}
\begin{aligned}
(x,y)&=\left(\sum_{i\in I}(x_i,x)x_i,\sum_{j\in I}(x_j,y)x_j\right)\\
&=\sum_{i,j\in I}(x,x_i)(x_j,y)\delta_{ij}\\
&=\sum_{i\in I}(x,x_i)(x_i,y).
\end{aligned}
\end{equation}
求和与内积的交换可以按通常方式论证: 先考察有限和, 再取极限.

(ii)$\Rightarrow$(iii): 令 $y=x$.

(iii)$\Rightarrow$(iv): 若 $x$ 与每个 $x_i$ 正交, 则 \eqref{eq:hilbert-basis-parseval-characterization} 意味着 $x=0$.

(iv)$\Rightarrow$(i): 令
\begin{equation}\label{eq:orthonormal-family-expansion-subspace}
Y:=\left\{x\in H\mid x=\sum_{i\in I}(x_i,x)x_i\right\}.
\end{equation}
设 $x^{(n)}$ 为 $Y$ 中的 Cauchy 序列. 那么存在至多可数多个 $i\in I$, 使 $(x_i,x^{(n)})\ne0$ 对任意 $n$ 成立. 用 $\widetilde I$ 表示这个至多可数的集合, 并令
\begin{equation}\label{eq:countable-orthonormal-expansion-subspace}
\widetilde Y:=\left\{x\in H\mid x=\sum_{i\in\widetilde I}(x_i,x)x_i\right\}.
\end{equation}
Parseval 等式表明 $\widetilde Y$ 或者有限维, 或者与序列空间 $\ell^2$ 等距, 因而完备; 见 Example~\ref{ex:ell-two-hilbert-space}. 因此, Cauchy 序列 $x^{(n)}$ 在 $\widetilde Y\subseteq Y$ 中有极限, 即 $Y$ 为 $H$ 的闭子空间. 另一方面, (iv) 表明 $Y^\perp=\{0\}$, 由 Theorem~\ref{thm:hilbert-orthogonal-projection} 得出 $Y=H$.
\end{Proof}

\begin{Corollary}
\label{cor:hilbert-basis-existence}
每个 Hilbert 空间都有一组基.
\end{Corollary}

\begin{Proof}
标准地应用 Zorn 引理可知, 存在极大标准正交集.
\end{Proof}

对可分 Hilbert 空间, 可以用 Schmidt 正交化过程以更具构造性的方式找到基. 设 $\{x_n\}_{n\in\mathbb{N}}$ 为可数稠密集. 从这个序列中删去每个能够表示为前面各元素的线性组合的元素. 这样得到一个新的线性无关元素序列 $\{y_n\}$, 且 $y_n$ 的线性张成空间仍在 $H$ 中稠密. 现在如下构造序列 $z_n$:
\begin{equation}\label{eq:schmidt-orthogonalization-recursion}
\begin{aligned}
z_1&=y_1/\|y_1\|,\\
u_n&=y_n-\sum_{i=1}^{n-1}(z_i,y_n)z_i,\\
z_n&=u_n/\|u_n\|.
\end{aligned}
\end{equation}
容易看出, $z_n$ 标准正交, 且其线性张成空间与 $y_n$ 的相同, 因而在 $H$ 中稠密. 因此 Theorem~\ref{thm:hilbert-basis-equivalent-characterizations} 的条件 (iv) 适用, $z_n$ 构成一组基.
